\renewcommand{\thetheorem}{7.\arabic{theorem}}
\setcounter{theorem}{0}
\setcounter{chapter}{6}
% ---------------------------------------------------------------------------
% Provenance of the prose in this file.
%   % Extractor: <model>   the block below was transcribed from Prof. Biran's
%                          handwritten notes by that model
%   % Creator:   <model>   the block below was written by that model and has no
%                          counterpart in the notes
% A tag holds until the next tag.
%
% This chapter was written before the convention existed, so the prose carries
% no tags except on the blocks added by later passes. Untagged prose is the
% original transcription.
% ---------------------------------------------------------------------------

\chapter{Vector Spaces}

\subsection{Definition and Axioms}

\begin{definition}[Vector Space]
\label{def:vector_space}
A \newterm{vector space} over a field $K$ is a set $V$, endowed with two operations:
\begin{align*}
+: V \times V &\longrightarrow V, \quad (v_1, v_2) \mapsto v_1 + v_2  &\text{(Addition of vectors)} \\
\cdot: K \times V &\longrightarrow V, \quad (\alpha, v) \mapsto \alpha \cdot v  &\text{(Scalar multiplication)}
\end{align*}
such that the following conditions (axioms) hold:
\begin{enumerate}[label=(\textbf{V}\arabic*), leftmargin=3.5em]
    \item $\forall v_1, v_2, v_3 \in V$: $v_1 + (v_2 + v_3) = (v_1 + v_2) + v_3$.
    \item $\exists$ an element $0 \in V$ such that $0 + v = v + 0 = v$ for all $v \in V$. (Sometimes denoted $0_V$).
    \item $\forall v \in V, \exists v' \in V$ such that $v + v' = 0$.
    \item $\forall v_1, v_2 \in V$: $v_1 + v_2 = v_2 + v_1$.
    \item $\forall a, b \in K, \forall v \in V$: $a \cdot (b \cdot v) = (a \cdot b) \cdot v$.
    \item $\forall v \in V$: $1 \cdot v = v$.
    \item $\forall \alpha \in K, \forall v_1, v_2 \in V$: $\alpha \cdot (v_1 + v_2) = \alpha \cdot v_1 + \alpha \cdot v_2$.
    \item $\forall \alpha_1, \alpha_2 \in K, \forall v \in V$: $(\alpha_1 + \alpha_2) \cdot v = \alpha_1 \cdot v + \alpha_2 \cdot v$.
\end{enumerate}
The elements of $K$ are called \newterm{scalars} and the elements of $V$ are called \newterm{vectors}.
\end{definition}

\begin{lemma}[Elementary Consequences of the Axioms]
\label{lem:vector_space_basics}
Let $V$ be a vector space over a field $K$. Then the following properties hold:
\begin{enumerate}[label=\textbf{(\alph*)}]
    \item The $0$ vector from (V2) is unique. To avoid confusion (e.g., with $0 \in K$), we will sometimes denote it by $0_V$.
    \item The element $v'$ (corresponding to $v$) from (V3) is unique. We will denote it by $-v$. We will also write $v - w := v + (-w)$ for all $v, w \in V$.
    \item For all $v \in V$, we have $0 \cdot v = 0_V$ (i.e., $0_K \cdot v = 0_V$).
    \item For all $\alpha \in K$, we have $\alpha \cdot 0_V = 0_V$.
    \item $(-1) \cdot v = -v$ for all $v \in V$.
    \item $-(-v) = v$ for all $v \in V$.
    \item If $a \cdot v = 0_V$ for some $a \in K$ and $v \in V$, then either $a = 0$ or $v = 0_V$.
    \item Associativity for $+$ holds also for sums of $n$ elements. That is, $v_1 + \dots + v_n$ makes sense no matter how we place the brackets:
    \[
    v_1 + \bigl(v_2 + (\dots + (v_{n-1} + v_n) \dots)\bigr),
    \qquad
    (v_1 + v_2) + (\dots), \qquad \text{and so on},
    \]
    all give the same element of $V$.
\end{enumerate}
\end{lemma}

\begin{proof}
\begin{enumerate}[label=\textbf{(\alph*)}]
    \item Assume there are two such zero elements, $0_1$ and $0_2$.
Then, by applying (V2) to both, we get: $0_1 = 0_1 + 0_2 = 0_2 + 0_1 = 0_2$. Thus, $0_1 = 0_2$.

    \item Assume there are two inverses $v_1'$ and $v_2'$ for an element $v$.
Then we can write:
\[
v_1' = v_1' + 0_V = v_1' + (v + v_2') = (v_1' + v) + v_2' = 0_V + v_2' = v_2'
\]
Thus, the inverse is unique.

    \item Consider $0 \cdot v = (0 + 0) \cdot v = 0 \cdot v + 0 \cdot v$, where the second equality is (V8).
By adding $-(0 \cdot v)$ to both sides, we yield $0_V = 0 \cdot v$.

    \item Similarly, $\alpha \cdot 0_V = \alpha \cdot (0_V + 0_V) = \alpha \cdot 0_V + \alpha \cdot 0_V$ by (V7). Adding $-(\alpha \cdot 0_V)$ yields $0_V = \alpha \cdot 0_V$.

    \item Observe that $v + (-1) \cdot v = 1 \cdot v + (-1) \cdot v = (1 - 1) \cdot v = 0 \cdot v = 0_V$. Since the inverse is unique by part \textbf{(b)}, it must be that $(-1) \cdot v = -v$.

    \item By definition, $-(-v)$ is the unique element that added to $-v$ gives $0_V$. But $v$ itself satisfies $(-v) + v = v + (-v) = 0_V$ by (V4) and (V3), so uniqueness in part \textbf{(b)} forces $-(-v) = v$.

    \item Let $a \in K$ and $v \in V$ and assume that $a \cdot v = 0_V$. If $a = 0$ there is nothing to prove, so suppose $a \ne 0$. Because $K$ is a field, there is an element $a^{-1} \in K$ with $a^{-1} \cdot a = 1$. Therefore,
    \[
    v \overset{\text{(V6)}}{=} 1 \cdot v = (a^{-1} \cdot a) \cdot v \overset{\text{(V5)}}{=} a^{-1} \cdot (a \cdot v) = a^{-1} \cdot 0_V \overset{\text{\textbf{(d)}}}{=} 0_V.
    \]
    So either $a = 0$ or $v = 0_V$, as claimed.

    \item This follows from (V1) by induction on $n$, exactly as for the associativity of addition in a field. From now on we will therefore write $v_1 + \dots + v_n$ without brackets.
\end{enumerate}
\end{proof}

\begin{exercise}[Non-Vector Space Structure on the Extended Reals]
\label{exc:extended_reals_not_vector_space}
Let $\infty$ and $-\infty$ denote two distinct formal elements outside $\mathbb{R}$. Define an addition and a scalar multiplication on the set $V := \mathbb{R} \cup \{\infty\} \cup \{-\infty\}$ by standard arithmetic on $\mathbb{R}$, and for $t \in \mathbb{R}$:
\[
t \cdot \infty := \begin{cases} -\infty & \text{if } t < 0, \\ 0 & \text{if } t = 0, \\ \infty & \text{if } t > 0, \end{cases}
\qquad
t \cdot (-\infty) := \begin{cases} \infty & \text{if } t < 0, \\ 0 & \text{if } t = 0, \\ -\infty & \text{if } t > 0, \end{cases}
\]
along with $t + \infty = \infty + t := \infty$, $t + (-\infty) = (-\infty) + t := -\infty$, $\infty + \infty := \infty$, $(-\infty) + (-\infty) := -\infty$, and $\infty + (-\infty) = (-\infty) + \infty := 0$.
Determine whether $V$ endowed with these operations forms a vector space over $\mathbb{R}$.
\exinfo{This exercise is Problem 4 of Problem Sheet 5, Linear Algebra I, Autumn Semester 2025.}
\end{exercise}

\subsection{The Coordinate Space \texorpdfstring{$K^n$}{K to the n}}

\begin{example}[The Coordinate Space \texorpdfstring{$K^n$}{K to the n}]
\label{ex:coordinate_space}
Let $K$ be a field and $n \in \mathbb{N}$. Put
\[
K^n := \{(a_1, \dots, a_n) \mid a_i \in K, \; 1 \le i \le n\}.
\]
The set $K^n$ becomes a vector space over $K$ if we endow it with the following operations.
Let $v = (a_1, \dots, a_n)$ and $w = (b_1, \dots, b_n) \in K^n$. We define addition in $K^n$ component-wise, using the addition of the field $K$ in each slot:
\[
v + w = (a_1, \dots, a_n) + (b_1, \dots, b_n) := (a_1 + b_1, \dots, a_n + b_n)
\]
For $\alpha \in K$, we define scalar multiplication as:
\[
\alpha \cdot v = \alpha \cdot (a_1, \dots, a_n) := (\alpha \cdot a_1, \dots, \alpha \cdot a_n)
\]
The zero vector is explicitly defined as $0 := (0, \dots, 0) \in K^n$.
With these operations $K^n$ becomes a vector space over $K$; the proof rests entirely on the fact that $K$ itself satisfies the field axioms. We call $K^n$ also the \newterm{$n$-dimensional coordinate space} over $K$.
\end{example}

\begin{exercise}
\label{exc:kn_axioms}
Check the details: verify each of the eight axioms (V1)--(V8) for $K^n$.
\end{exercise}

\begin{exercise}[Multiplicative Vector Space Structure on Positive Real Pairs]
\label{exc:positive_reals_vector_space_mc}
Consider the set of pairs of strictly positive real numbers $\mathbb{R}_{>0}^2 := \{(x_1, x_2) \in \mathbb{R}^2 \mid x_1 > 0, \, x_2 > 0\}$. Define an addition operation on $\mathbb{R}_{>0}^2$ by componentwise multiplication:
\[
\begin{pmatrix} x_1 \\ x_2 \end{pmatrix} \oplus \begin{pmatrix} y_1 \\ y_2 \end{pmatrix} := \begin{pmatrix} x_1 y_1 \\ x_2 y_2 \end{pmatrix}.
\]
Consider the following three candidate definitions of scalar multiplication for $\lambda \in \mathbb{R}$:
\begin{enumerate}[label=\textbf{(\arabic*)}]
    \item $\lambda \odot \begin{pmatrix} x_1 \\ x_2 \end{pmatrix} := \begin{pmatrix} \lambda x_1 \\ \lambda x_2 \end{pmatrix}$;
    \item $\lambda \odot \begin{pmatrix} x_1 \\ x_2 \end{pmatrix} := \begin{pmatrix} e^{\lambda} x_1 \\ e^{\lambda} x_2 \end{pmatrix}$;
    \item $\lambda \odot \begin{pmatrix} x_1 \\ x_2 \end{pmatrix} := \begin{pmatrix} x_1^{\lambda} \\ x_2^{\lambda} \end{pmatrix}$.
\end{enumerate}
Under which of these candidate definitions does $(\mathbb{R}_{>0}^2, \oplus, \odot)$ form a valid vector space over $\mathbb{R}$?
\exinfo{This exercise is Multiple Choice Question 2 of Problem Sheet 5, Linear Algebra I, Autumn Semester 2025.}
\end{exercise}

\begin{importantremark}
\label{rem:row_or_column}
Sometimes we write the elements of $K^n$ as row vectors and sometimes as column vectors, whichever is more convenient at the time. The two conventions describe the same vector space, and we will switch between them without further comment, but whenever a matrix product $A \cdot x$ is involved, the elements of $K^n$ are always to be read as column vectors, as in \cref{def:matrix_vector_multiplication}.
\end{importantremark}

\subsection{Linear Subspaces}

\begin{definition}[Linear Subspace]
\label{def:linear_subspace}
Let $V$ be a vector space over $K$. A subset $U \subseteq V$ is called a \newterm{linear subspace} of $V$ (or just a \newterm{subspace} of $V$) if the following holds:
\begin{enumerate}[label=(\textbf{LSS}\arabic*), leftmargin=4.5em]
    \item $U \ne \emptyset$.
    \item For all $u_1, u_2 \in U$, we have $u_1 + u_2 \in U$ (closed under addition).
    \item For all $\alpha \in K$ and all $u \in U$, we have $\alpha \cdot u \in U$ (closed under scalar multiplication).
\end{enumerate}
\end{definition}

\begin{lemma}[A Subspace is a Vector Space]
\label{lem:subspace_is_vector_space}
Let $V$ be a vector space over $K$ and let $U \subseteq V$ be a subspace. Then $U$ is a vector space in its own right, when endowed with the operations inherited from $V$.
\end{lemma}

\begin{proof}
Conditions (LSS2) and (LSS3) say precisely that the operations $+$ and $\cdot$ of $V$ restrict to operations
\[
+ \colon U \times U \longrightarrow U, \qquad \cdot \colon K \times U \longrightarrow U,
\]
so that $U$ carries the two operations required by \cref{def:vector_space}. That these operations satisfy the eight axioms of a vector space then follows by direct verification, since each axiom is an identity that already holds for all elements of $V$ and therefore in particular for all elements of $U$.

The one axiom that deserves a word is (V2), because it asserts the \emph{existence} of an element of $U$ rather than an identity. Pick any element $u \in U$; this is possible because $U \ne \emptyset$ by (LSS1). By (LSS3),
\[
0_V = 0_K \cdot u \in U,
\]
where the first equality is \cref{lem:vector_space_basics} \textbf{(c)}. Now $0_V + w = w$ holds for every $w \in U$, because it already holds for every $w \in V$ by (V2) for $V$. Hence $0_V$ serves as the zero element of $U$.
\end{proof}

\begin{exercise}
\label{exc:subspace_axioms}
Carry out the direct verification of the remaining seven axioms for $U$.
\end{exercise}

In practice one rarely checks (LSS1)--(LSS3) separately. The following criterion compresses them into two conditions, one of which is a single membership test.

\begin{lemma}[Criterion for a Linear Subspace]
\label{lem:subspace_criterion}
Let $V$ be a vector space over $K$, and let $U \subseteq V$ be a subset. Then $U$ is a linear subspace of $V$ if and only if the following holds:
\begin{enumerate}[label=\textbf{(\alph*)}]
    \item $0_V \in U$.
    \item For all $\alpha_1, \alpha_2 \in K$ and all $u_1, u_2 \in U$, we have $\alpha_1 u_1 + \alpha_2 u_2 \in U$.
\end{enumerate}
\end{lemma}

\begin{proof}
\textbf{Direction 1 ($\implies$):} Suppose $U \subseteq V$ is a subspace. Then $0_V \in U$ was established in the proof of \cref{lem:subspace_is_vector_space}, which gives \textbf{(a)}. For \textbf{(b)}, given $\alpha_1, \alpha_2 \in K$ and $u_1, u_2 \in U$, condition (LSS3) gives $\alpha_1 u_1 \in U$ and $\alpha_2 u_2 \in U$, and then (LSS2) gives $\alpha_1 u_1 + \alpha_2 u_2 \in U$.

\textbf{Direction 2 ($\impliedby$):} Suppose \textbf{(a)} and \textbf{(b)} hold. Then $U \ne \emptyset$ by \textbf{(a)}, which is (LSS1). Let $u_1, u_2 \in U$. By \textbf{(b)} applied with $\alpha_1 = \alpha_2 = 1$,
\[
\underbrace{1 \cdot u_1 + 1 \cdot u_2}_{= \; u_1 + u_2} \in U,
\]
which is (LSS2). Finally, let $\alpha \in K$ and $u \in U$. By \textbf{(b)} applied with $\alpha_1 = \alpha$, $\alpha_2 = 0$ and $u_1 = u_2 = u$,
\[
\underbrace{\alpha \cdot u + 0 \cdot u}_{= \; \alpha \cdot u} \in U,
\]
which is (LSS3). Hence $U$ is a linear subspace.
\end{proof}

\begin{exercise}[Verifying Linear Subspaces over \texorpdfstring{$\mathbb{R}$}{R} and \texorpdfstring{$\mathbb{F}_2$}{F2}]
\label{exc:verifying_subspaces_r_f2}
Determine which of the following subsets are linear subspaces of the given vector spaces:
\begin{enumerate}
    \item $S_1 := \{(x_1, x_2, x_3) \in \mathbb{R}^3 \mid x_1 = x_2 = 2x_3\} \subseteq \mathbb{R}^3$.
    \item $S_2 := \{(x_1, x_2) \in \mathbb{R}^2 \mid x_1^2 + x_2^4 = 0\} \subseteq \mathbb{R}^2$.
    \item $S_3 := \{(\mu + \lambda, \lambda^2) \in \mathbb{R}^2 \mid \mu, \lambda \in \mathbb{R}\} \subseteq \mathbb{R}^2$.
\end{enumerate}
What changes in parts (b) and (c) if the base field $\mathbb{R}$ is replaced by the finite field $\mathbb{F}_2$?
\exinfo{This exercise is Problem 2 of Problem Sheet 5, Linear Algebra I, Autumn Semester 2025.}
\end{exercise}

\begin{exercise}[Subspace Identification in Coordinate Spaces]
\label{exc:subspace_identification_mc}
Which of the following subsets are linear subspaces of the ambient coordinate spaces?
\begin{enumerate}
    \item $\{(x_1, x_2, x_3) \in \mathbb{R}^3 \mid 3x_1 + 5x_2 + 3x_3 = 0 \text{ and } 2x_2 + x_3 = 0\} \subseteq \mathbb{R}^3$.
    \item $\{(x_1, x_2, x_3) \in \mathbb{R}^3 \mid x_1 + x_2 + x_3 = 3\} \subseteq \mathbb{R}^3$.
    \item $\{(x_1, x_2) \in \mathbb{R}^2 \mid x_1 > x_2\} \subseteq \mathbb{R}^2$.
    \item $\{(0, x, 2x, 3x) \mid x \in \mathbb{R}\} \subseteq \mathbb{R}^4$.
    \item $\{(x^4, x^3, x^2, x) \mid x \in \mathbb{R}\} \subseteq \mathbb{R}^4$.
\end{enumerate}
\exinfo{This exercise is Multiple Choice Question 1 of Problem Sheet 5, Linear Algebra I, Autumn Semester 2025.}
\end{exercise}

\subsection{Examples of Subspaces}

\begin{example}[The Trivial Subspace]
\label{ex:trivial_subspace}
Let $V$ be a vector space over $K$. Then $\{0_V\} \subseteq V$ is a linear subspace: it is non-empty, and $0_V + 0_V = 0_V$ as well as $\alpha \cdot 0_V = 0_V$ for every $\alpha \in K$, by \cref{lem:vector_space_basics} \textbf{(d)}. At the other extreme, $V \subseteq V$ is of course a subspace of itself.
\end{example}

\begin{example}[A Plane Through the Origin, and One That Is Not]
\label{ex:plane_subspace}
Fix $b \in K$ and consider the subset
\[
U_b := \{(x_1, x_2, x_3) \in K^3 \mid x_1 + x_2 + x_3 = b\} \subseteq K^3.
\]

\begin{claim*}[$U_b$ is a subspace exactly for $b = 0$]
\label{claim:U_b_subspace}
$U_b \subseteq K^3$ is a linear subspace if and only if $b = 0$.
\end{claim*}

\emph{Proof of the direction $\implies$.} Suppose $U_b$ is a subspace of $K^3$. Let $(x_1, x_2, x_3)$ and $(y_1, y_2, y_3) \in U_b$. Then by definition,
\begin{equation} \label{eq:U_b_membership}
x_1 + x_2 + x_3 = b, \qquad y_1 + y_2 + y_3 = b.
\end{equation}
Since $U_b$ is a subspace, it is closed under addition, and therefore, $(x_1 + y_1, x_2 + y_2, x_3 + y_3) \in U_b$. By the definition of $U_b$, this implies
\[
(x_1 + y_1) + (x_2 + y_2) + (x_3 + y_3) = b.
\]
But rearranging the terms and using \eqref{eq:U_b_membership}, we also get
\[
(x_1 + y_1) + (x_2 + y_2) + (x_3 + y_3) = (x_1 + x_2 + x_3) + (y_1 + y_2 + y_3) = b + b.
\]
Therefore, $b + b = b$, which implies $b = 0$. \hfill $\diamond$
\end{example}

\begin{exercise}
\label{exc:U_b_converse}
Prove the direction $\impliedby$ of the claim on \cpageref{claim:U_b_subspace} that $U_b$ is a subspace exactly for $b = 0$: show that if $b = 0$, then $U_0$ satisfies (LSS1), (LSS2) and (LSS3), and is therefore a linear subspace of $K^3$.
\end{exercise}

\begin{exercise}[Unions of Three Linear Subspaces]
\label{exc:union_three_subspaces}
Let $V$ be a vector space over a field $K$, and let $V_1, V_2, V_3 \subseteq V$ be linear subspaces such that none of the three is contained in any other.
Determine with proof whether the union $V_1 \cup V_2 \cup V_3$ is always, sometimes, or never a linear subspace of $V$.
\exhint{Consider different choices for the base field $K$, such as $\mathbb{F}_2$.}
\exinfo{This exercise is Problem 6 of Problem Sheet 5, Linear Algebra I, Autumn Semester 2025.}
\end{exercise}

\subsection{The Solution Set of a Linear System}

The previous example is a special case of something much more important. It answers the third of the questions we set ourselves in the previous chapter, in the list of goals on \cpageref{goals:solution_set}: the solution set of a system of linear equations does carry a structure of its own, provided the system is homogeneous.

Let $A \in \M_{m \times n}(K)$ and fix $b \in K^m$, viewed as a column vector. Consider
\[
L := \{x \in K^n \mid A \cdot x = b\} \subseteq K^n,
\]
where the elements of $K^n$ are viewed as column vectors as well, following the remark on row and column notation on \cpageref{rem:row_or_column}. This is exactly the set $L(S)$ of solutions of the system $(S): Ax = b$ from \eqref{eq:matrix_equation}.

\begin{lemma}[Linearity of Matrix-Vector Multiplication]
\label{lem:matrix_vector_linearity}
Let $A \in \M_{m \times n}(K)$. For all $u, v \in K^n$ and all $a \in K$ we have:
\begin{enumerate}[label=\textbf{(\alph*)}]
    \item $A \cdot (u + v) = A \cdot u + A \cdot v$.
    \item $A \cdot (a \cdot u) = a \cdot (A \cdot u)$.
\end{enumerate}
\end{lemma}

\begin{proof}
Write $A = (a_{ij})$ with $1 \le i \le m$, $1 \le j \le n$, and let
\[
u = \begin{pmatrix} u_1 \\ \vdots \\ u_n \end{pmatrix}, \qquad
v = \begin{pmatrix} v_1 \\ \vdots \\ v_n \end{pmatrix}.
\]
We prove \textbf{(a)} by comparing the $i$-th entries of both sides. By \cref{def:matrix_vector_multiplication}, entry number $i$ of $A \cdot u$ is $a_{i1}u_1 + \dots + a_{in}u_n$, and entry number $i$ of $A \cdot v$ is $a_{i1}v_1 + \dots + a_{in}v_n$. Adding these in $K$ and regrouping by distributivity gives
\begin{align*}
\bigl(A \cdot u + A \cdot v\bigr)_i
&= \bigl(a_{i1}u_1 + \dots + a_{in}u_n\bigr) + \bigl(a_{i1}v_1 + \dots + a_{in}v_n\bigr) \\
&= a_{i1}(u_1 + v_1) + a_{i2}(u_2 + v_2) + \dots + a_{in}(u_n + v_n) \\
&= \bigl(A \cdot (u + v)\bigr)_i,
\end{align*}
the last equality again by \cref{def:matrix_vector_multiplication}, since the $j$-th entry of $u + v$ is $u_j + v_j$. As this holds for every $1 \le i \le m$, the two column vectors are equal.
\end{proof}

\begin{exercise}
\label{exc:matrix_vector_scalar}
Prove statement \textbf{(b)} of \cref{lem:matrix_vector_linearity}.
\end{exercise}

\begin{theorem}[When the Solution Set is a Subspace]
\label{thm:solution_set_subspace}
$L \subseteq K^n$ is a linear subspace if and only if $b = 0 \in K^m$.
\end{theorem}

\begin{proof}
\textbf{Direction 1 ($\impliedby$):} We assume $b = 0$ and have to show that $L \subseteq K^n$ is a subspace.

First, $x = 0 \in L$, since indeed $A \cdot 0 = 0$. Next, if $x = (x_1, \dots, x_n) \in L$ and $y = (y_1, \dots, y_n) \in L$, then by \cref{lem:matrix_vector_linearity} \textbf{(a)},
\[
A \cdot (x + y) = A \cdot x + A \cdot y = 0 + 0 = 0,
\]
and therefore, $x + y \in L$. This shows that $x, y \in L$ implies $x + y \in L$. Finally, let $x \in L$ and $a \in K$. By \cref{lem:matrix_vector_linearity} \textbf{(b)},
\[
A \cdot (a \cdot x) = a \cdot (A \cdot x) = a \cdot 0 = 0,
\]
so $a \cdot x \in L$. This shows $L$ is a subspace.

\textbf{Direction 2 ($\implies$):} Suppose $L \subseteq K^n$ is a subspace. We need to show that $b = 0$. Since $L$ is a subspace, $0 \in L$ by \cref{lem:subspace_criterion} \textbf{(a)}. By the definition of $L$, this means $A \cdot 0 = b$. But $A \cdot 0 = 0$, and therefore, $b = 0$.
\end{proof}

\begin{importantremark}
\label{rem:homogeneous_vs_inhomogeneous}
So the solution set of a \emph{homogeneous} system is a linear subspace of $K^n$, while the solution set of an inhomogeneous one is never a subspace; it does not even contain $0$. It is not, however, shapeless: recall the last example of the previous chapter, whose general solution had the form \qt{one particular solution plus an arbitrary multiple of a fixed vector}. We will make that observation precise once we have the notion of dimension.
\end{importantremark}

\begin{exercise}[Parametric Linear Systems and Subspaces]
\label{exc:parametric_linear_system_subspace}
Let $m \in \mathbb{R}$. Consider the system of linear equations in $(x, y) \in \mathbb{R}^2$:
\[
\begin{cases}
\phantom{m}x + my = -3 \\
mx + 4y = \phantom{-}6
\end{cases}
\]
\begin{enumerate}
    \item Determine the set of solutions $S \subseteq \mathbb{R}^2$ depending on the parameter $m$.
    \item For which values of $m$ is $S$ a linear subspace of $\mathbb{R}^2$?
    \item Provide a geometric interpretation of $S$ in the real plane $\mathbb{R}^2$ for each case of $m$.
\end{enumerate}
\exinfo{This exercise is Problem 1 of Problem Sheet 5, Linear Algebra I, Autumn Semester 2025.}
\end{exercise}

\subsection{More Examples of Vector Spaces}

\subsubsection{Spaces of Sequences \texorpdfstring{$K^\infty$}{K to the infinity}}

\begin{example}[Spaces of Sequences]
\label{ex:sequence_space}
Let $K$ be a field. Define
\[
K^\infty := \{(a_0, a_1, a_2, \dots, a_k, \dots) \mid a_i \in K \; \forall i\},
\]
the set of all infinite sequences with entries in $K$. We define $+$ and $\cdot$ on $K^\infty$ entry-wise:
\begin{align*}
(a_0, a_1, \dots, a_k, \dots) + (b_0, b_1, \dots, b_k, \dots) &:= (a_0 + b_0, a_1 + b_1, \dots, a_k + b_k, \dots), \\
a \cdot (a_0, a_1, \dots, a_k, \dots) &:= (a \cdot a_0, a \cdot a_1, \dots, a \cdot a_k, \dots) \quad \text{for } a \in K.
\end{align*}
\end{example}

\begin{exercise}
\label{exc:sequence_space_axioms}
With these operations, $K^\infty$ is a vector space over $K$.
\end{exercise}

\begin{example}[Fibonacci Sequences as a Subspace]
\label{ex:fib_subspace}
Recall the Fibonacci sequences from \cref{prop:fib_closure}:
\[
\Fib = \{(a_0, a_1, \dots, a_k, \dots) \in \mathbb{R}^\infty \mid a_k = a_{k-1} + a_{k-2} \; \forall k \ge 2\}.
\]
Then $\Fib \subseteq \mathbb{R}^\infty$ is a linear subspace, which is precisely what we verified by hand in the very first chapter, before we had the word for it.

More generally, if we fix $\alpha, \beta \in \mathbb{R}$, then
\[
U_{\alpha, \beta} := \{(a_0, a_1, \dots, a_k, \dots) \in \mathbb{R}^\infty \mid a_k = \alpha a_{k-1} + \beta a_{k-2} \; \forall k \ge 2\}
\]
is also a linear subspace of $\mathbb{R}^\infty$, with $\Fib = U_{1,1}$.
\end{example}

\subsubsection{Spaces of Functions \texorpdfstring{$K^S$}{K to the S}}

\begin{example}[Function Spaces]
\label{ex:function_space}
Fix a field $K$ and a non-empty set $S$. We define the function space
\[
K^S := \{f \colon S \to K \mid f \text{ is a function}\}.
\]
We define addition $+$ and scalar multiplication $\cdot$ on $K^S$ as follows. Let $f, g \in K^S$ and $a \in K$. We define $(f + g) \colon S \to K$ point-wise by $x \mapsto f(x) + g(x)$, where the addition on the right is that of $K$, and we define $(a \cdot f) \colon S \to K$ point-wise by $x \mapsto a \cdot f(x)$.
\end{example}

\begin{exercise}
\label{exc:function_space_axioms}
$K^S$ is a vector space over $K$.
\end{exercise}

\begin{exercise}[Even and Odd Functions as Linear Subspaces]
\label{exc:even_odd_function_subspaces}
Let $K$ be a field with $1 + 1 \ne 0$, and consider the function space $V := K^K = \mathrm{Abb}(K, K)$ endowed with pointwise addition and scalar multiplication. Define the subsets:
\begin{align*}
V_{\mathrm{even}} &:= \{f \colon K \to K \mid f(-x) = f(x) \;\forall x \in K\}, \\
V_{\mathrm{odd}} &:= \{f \colon K \to K \mid f(-x) = -f(x) \;\forall x \in K\}.
\end{align*}
\begin{enumerate}
    \item Show that $V_{\mathrm{even}}$ and $V_{\mathrm{odd}}$ are linear subspaces of $V$.
    \item Show that $V_{\mathrm{even}} + V_{\mathrm{odd}} = V$, where $V_{\mathrm{even}} + V_{\mathrm{odd}} := \{v + w \mid v \in V_{\mathrm{even}}, w \in V_{\mathrm{odd}}\}$.
    \item Show that $V_{\mathrm{even}} \cap V_{\mathrm{odd}} = \{0\}$.
\end{enumerate}
\exinfo{This exercise is Problem 3 of Problem Sheet 5, Linear Algebra I, Autumn Semester 2025.}
\end{exercise}

\begin{exercise}[The Power Set as an \texorpdfstring{$\mathbb{F}_2$}{F2}-Vector Space]
\label{exc:power_set_f2_vector_space}
Let $X$ be a set and let $\mathcal{P}(X)$ be its power set. For all $A, B \in \mathcal{P}(X)$ and $\lambda \in \mathbb{F}_2 = \{0, 1\}$, define the operations:
\begin{align*}
A \bigtriangleup B &:= (A \cup B) \setminus (A \cap B), \\
\lambda \cdot A &:= \begin{cases} \emptyset & \text{if } \lambda = 0, \\ A & \text{if } \lambda = 1. \end{cases}
\end{align*}
Show that $(\mathcal{P}(X), \bigtriangleup, \cdot, \emptyset)$ is a vector space over the field $\mathbb{F}_2$.
\exinfo{This exercise is Problem 5 of Problem Sheet 5, Linear Algebra I, Autumn Semester 2025.}
\end{exercise}

\begin{example}[Continuous Functions]
\label{ex:continuous_subspace}
Take $K = \mathbb{R}$ and let $S := [0,1] = \{x \in \mathbb{R} \mid 0 \le x \le 1\}$. Define
\[
C(S) := \{f \in \mathbb{R}^S \mid f \text{ is continuous on } [0,1]\} \subseteq \mathbb{R}^S.
\]
Then $C(S) \subseteq \mathbb{R}^S$ is a linear subspace.
\end{example}

\begin{exercise}
\label{exc:continuous_subspace}
Prove that $C(S)$ is a linear subspace of $\mathbb{R}^S$: check that the constant function $0$ is continuous, and that a sum of two continuous functions, as well as a scalar multiple of a continuous function, is again continuous on $[0,1]$.
\end{exercise}

\subsubsection{Polynomials \texorpdfstring{$K[x]$}{K[x]}}

\begin{example}[Polynomial Spaces]
\label{ex:polynomial_space}
Let $K$ be a field and let $x$ be a \emph{formal} variable; we may think of $x$ as a \qt{free element of $K$}. A \newterm{polynomial} over $K$ (or polynomial with coefficients in $K$) is an expression of the form
\[
f(x) = a_0 + a_1 x + a_2 x^2 + \dots + a_n x^n, \qquad n \in \mathbb{N},
\]
where $a_i \in K$ for all $0 \le i \le n$. The $a_i$ are called the \newterm{coefficients} of $f$. The \newterm{zero polynomial} is the one with $a_0 = a_1 = \dots = 0$. Two polynomials are equal if and only if all their corresponding coefficients are equal.

To add two polynomials
\[
f(x) = a_0 + a_1 x + \dots + a_n x^n, \qquad g(x) = b_0 + b_1 x + \dots + b_m x^m,
\]
we first pad the shorter one with zeros: if $n < m$ we set $a_{n+1} = \dots = a_m = 0$, and if $n > m$ we set $b_{m+1} = \dots = b_n = 0$. Putting $r := \max\{n, m\}$, both are then written with the same number of terms and we may define
\[
f(x) + g(x) := (a_0 + b_0) + (a_1 + b_1)x + (a_2 + b_2)x^2 + \dots + (a_r + b_r)x^r.
\]
If $c \in K$, we define
\[
c \cdot f(x) := c \cdot a_0 + c \cdot a_1 x + c \cdot a_2 x^2 + \dots + c \cdot a_n x^n.
\]
\end{example}

\begin{notation}
\label{not:polynomial_ring}
We write $K[x]$ for the set of all polynomials over $K$.
\end{notation}

\begin{exercise}
\label{exc:polynomial_space_axioms}
With the operations above, $K[x]$ is a vector space over $K$.
\end{exercise}

If we view polynomials as functions $K \to K$, then the polynomials over $K$ give us a linear subspace of $K^K$. This second point of view is however genuinely coarser than the first, and the following warning explains why.

\begin{importantremark}
\label{rem:polynomial_vs_function}
For some fields (e.g., $K = \mathbb{F}_2$, or $\mathbb{F}_p$ with $p$ prime) two \emph{different} polynomials may give the same function. For example, over $K = \mathbb{F}_2$ the polynomials $f(x) = x$ and $g(x) = x^2$ produce the same output value for every input in $\mathbb{F}_2$. But \textbf{as polynomials} $f(x)$ and $g(x)$ are strictly different objects, since their coefficients differ.
\end{importantremark}

\subsubsection{The Degree of a Polynomial}

Let $f(x) = a_0 + a_1 x + a_2 x^2 + \dots + a_n x^n$ be a polynomial with $a_n \ne 0$. We say that the \newterm{degree} of $f$ is $n$, and we write $\deg(f) := n \in \mathbb{Z}_{\ge 0}$. If $a_0 = a_1 = \dots = a_n = 0$, meaning $f$ is the zero polynomial, we set $\deg(f) := -\infty$. Consequently, $\deg(f) = -\infty$ if and only if $f = 0$.

\begin{remark}
Every polynomial has a well-defined degree: if $f \ne 0$ and $f(x) = a_0 + a_1 x + \dots + a_r x^r$, then
\[
\deg(f) := \max \{k \in \mathbb{Z}_{\ge 0} \mid a_k \ne 0\},
\]
and this maximum is taken over a non-empty finite set, hence exists.
\end{remark}

\begin{example}[Polynomials of Bounded Degree]
\label{ex:bounded_degree_subspace}
Fix $d \in \mathbb{Z}_{\ge 0} \cup \{-\infty\}$ and define
\[
K[x]_d := \{f \in K[x] \mid \deg(f) \le d\},
\]
where we agree that $-\infty \le k$ for all $k \in \mathbb{Z}$.

\begin{claim*}[{$K[x]_d$ is a subspace}]
\label{claim:bounded_degree_subspace}
$K[x]_d \subseteq K[x]$ is a linear subspace of $K[x]$.
\end{claim*}

Later, we will see that $K[x]_d$ is a subspace of dimension $d+1$.
\end{example}

\begin{question}
\label{q:degree_exactly_d}
\begin{enumerate}[label=\textbf{(\alph*)}]
    \item Is the set $\{f \in K[x] \mid \deg(f) = d\}$ also a subspace of $K[x]$?
    \item Consider the subset
    \[
    \left\{ p \in K[x]_5 \;\middle|\; p(x) = a_0 + a_1 x + \dots + a_5 x^5, \; \sum_{i=0}^{5} a_i = 0 \right\} \subseteq K[x]_5,
    \]
    that is, all polynomials of degree $\le 5$ for which the sum of the coefficients is $0$. Is this subset a linear subspace of $K[x]_5$?
\end{enumerate}
\end{question}

\subsubsection{Spaces of Matrices}

\begin{example}[Matrix Spaces]
\label{ex:matrix_space}
Consider the set $\M_{m \times n}(K)$ of all $m \times n$ matrices with entries in $K$. We define addition of matrices component-wise: if $A = (a_{ij})$ and $B = (b_{ij})$, then $A + B := (a_{ij} + b_{ij})$. Similarly, for $\alpha \in K$ we define scalar multiplication component-wise by $\alpha \cdot A := (\alpha \, a_{ij})$.

\begin{claim*}[$\M_{m \times n}(K)$ is a vector space]
\label{claim:matrix_space}
With these operations, $\M_{m \times n}(K)$ is a vector space over $K$. Its zero vector is the $m \times n$ matrix all of whose entries are $0$.
\end{claim*}
\end{example}

\begin{ainote}
The lecture notes index sequences from $1$, writing $K^\infty = \{(a_1, a_2, \dots)\}$ and
$\Fib$ with the recursion starting at $k \ge 3$. This transcript indexes from $0$
instead, so that $\Fib$ here is literally the same object as the one built in
\cref{prop:fib_closure}; nothing else changes. The notes also use $W$ for a
subspace, where this transcript writes $U$, matching the later chapters. The
proofs of \cref{lem:vector_space_basics} \textbf{(e)}, \textbf{(f)} and
\textbf{(h)} are left as exercises in the notes and are supplied here; the
exercises marked as such above are the ones the notes leave open.
\end{ainote}

% Creator: Opus 5
\subsection{Solutions to the Exercises}
\label{sec:solutions_vector_spaces}

\begin{exercisesolution}[Proof of \cref{exc:kn_axioms}, \cref{exc:sequence_space_axioms}, \cref{exc:function_space_axioms} and \cref{exc:polynomial_space_axioms}]
These four exercises ask for the same verification four times, and it is worth doing once in a form that covers all of them, because the reason they are true is the same reason each time.

In each case the underlying set is a set of $K$-valued functions on an index set $I$, and both operations are defined \emph{index by index}:
\begin{center}
\begin{tabular}{|l|l|l|l|}
\hline
space & index set $I$ & element & restriction \\
\hline
$K^n$ & $\{1, \dots, n\}$ & $(x_1, \dots, x_n)$ & none \\
$K^\infty$ & $\mathbb{Z}_{\ge 0}$ & $(a_0, a_1, \dots)$ & none \\
$K^S$ & $S$ & $f \colon S \to K$ & none \\
$K[x]$ & $\mathbb{Z}_{\ge 0}$ & coefficients $(a_0, a_1, \dots)$ & finitely many $a_k \neq 0$ \\
\hline
\end{tabular}
\end{center}
Write $u(i)$ for the entry of $u$ at index $i$. In every one of the four, the operations are
\[
    (u + v)(i) = u(i) + v(i), \qquad (a \cdot u)(i) = a \cdot u(i) \qquad \text{for all } i \in I,
\]
with the right-hand sides computed in $K$.

Now look at any one of the eight axioms, say (V1), $u + (v + w) = (u + v) + w$. Two elements of the space are equal exactly when they agree at every index, so the axiom holds if and only if
\[
    u(i) + \bigl(v(i) + w(i)\bigr) = \bigl(u(i) + v(i)\bigr) + w(i) \qquad \text{for every } i \in I,
\]
and that is the associativity axiom \textbf{(K1)} of the field $K$, applied once for each $i$. The same argument settles the other seven: each of (V1)--(V8) is an identity between two expressions built from $+$ and $\cdot$, both sides evaluate index by index, and the resulting identity in $K$ is one of \textbf{(K1)}--\textbf{(K10)}. Concretely, (V1) and (V3) come from \textbf{(K1)} and \textbf{(K2)}; the zero element is the function with $0$ at every index and the inverse of $u$ is $i \mapsto -u(i)$, so (V2) and (V4) come from \textbf{(K3)} and \textbf{(K4)}; (V5) and (V6) come from \textbf{(K5)} and \textbf{(K6)}; and (V7) and (V8) come from the two halves of distributivity, \textbf{(K8)}.

For $K[x]$ one extra word is needed, and it is the only difference between the four cases: one must know that the operations do not leave the set. If $u$ has $a_k = 0$ for all $k > N$ and $v$ has $b_k = 0$ for all $k > M$, then $u + v$ has all coefficients beyond $\max(N, M)$ equal to zero, and $a \cdot u$ has all coefficients beyond $N$ equal to zero. So both results are again polynomials, and the axioms then hold for the reason above. This finiteness condition is exactly what makes $K[x]$ smaller than $K^\infty$, and \cref{exc:k_infinity_independent} in ch. 9 turns on the same distinction.
\end{exercisesolution}

\begin{exercisesolution}[Proof of \cref{exc:subspace_axioms}]
Let $U \subseteq V$ be a linear subspace and let $u, v, w \in U$ and $a, b \in K$. All eight axioms are statements about elements of $U$, and every element of $U$ is an element of $V$, where the axioms already hold. For the seven axioms that are pure identities, (V1) and (V3) through (V8), there is therefore nothing to check beyond one point: that both sides of the identity are elements of $U$ at all. That is exactly what (LSS2) and (LSS3) of \cref{def:linear_subspace} provide, since each side is built from $u, v, w$ by additions and scalar multiplications.

Two of them do assert existence and so deserve a word. For (V2), the proof of \cref{lem:subspace_is_vector_space} above shows that the zero of $V$ lies in $U$ and serves as a zero there. For (V4), given $u \in U$ the element $(-1) \cdot u$ lies in $U$ by (LSS3), and it satisfies $u + (-1) \cdot u = 0_V$ by \cref{lem:vector_space_basics}, so it is an additive inverse of $u$ inside $U$.

The moral is worth keeping: a subspace inherits every identity for free, and the only things ever in question are closure and the presence of $0$. This is why \cref{lem:subspace_criterion} can compress the whole check into two conditions.
\end{exercisesolution}

\begin{exercisesolution}[Proof of \cref{exc:U_b_converse}]
Assume $b = 0$, so that $U_0 = \{(x_1, x_2, x_3) \in K^3 \mid x_1 + x_2 + x_3 = 0\}$. We verify the three conditions of \cref{def:linear_subspace}.

\textbf{(LSS1)} The zero vector $(0,0,0)$ satisfies $0 + 0 + 0 = 0$, so $0 \in U_0$ and $U_0 \neq \emptyset$.

\textbf{(LSS2)} Let $x, y \in U_0$, so $x_1 + x_2 + x_3 = 0$ and $y_1 + y_2 + y_3 = 0$. Then
\[
    (x_1 + y_1) + (x_2 + y_2) + (x_3 + y_3) = (x_1 + x_2 + x_3) + (y_1 + y_2 + y_3) = 0 + 0 = 0,
\]
using commutativity and associativity of addition in $K$ to regroup. Hence $x + y \in U_0$.

\textbf{(LSS3)} Let $x \in U_0$ and $a \in K$. Then
\[
    a x_1 + a x_2 + a x_3 = a \cdot (x_1 + x_2 + x_3) = a \cdot 0 = 0
\]
by distributivity \textbf{(K8)} and \cref{lem:mult_by_zero}. Hence $a \cdot x \in U_0$.

So $U_0$ is a linear subspace of $K^3$. Together with the direction proved in the text, this establishes the Claim on \cpageref{claim:U_b_subspace}. Note where the two halves differ: the forward direction used only closure under addition, and produced the equation $b + b = b$; the converse needs all three conditions, but each is a one-line consequence of the field axioms.
\end{exercisesolution}

\begin{exercisesolution}[Proof of \cref{exc:matrix_vector_scalar}]
Let $A = (a_{ij})$ with $1 \le i \le m$ and $1 \le j \le n$, let $u \in K^n$ be a column vector and let $a \in K$. We compare the $i$-th entries of the two sides. By \cref{def:matrix_vector_multiplication}, the $j$-th entry of $a \cdot u$ is $a \cdot u_j$, and therefore,
\[
    \bigl(A \cdot (a \cdot u)\bigr)_i = a_{i1}(a u_1) + a_{i2}(a u_2) + \dots + a_{in}(a u_n).
\]
Each summand may be rewritten as $a \cdot (a_{ij} u_j)$, using commutativity and associativity of multiplication in $K$, axioms \textbf{(K5)} and \textbf{(K10)}. Pulling the common factor $a$ out by distributivity \textbf{(K8)},
\[
    \bigl(A \cdot (a \cdot u)\bigr)_i = a \cdot \bigl(a_{i1}u_1 + \dots + a_{in}u_n\bigr) = a \cdot \bigl(A \cdot u\bigr)_i = \bigl(a \cdot (A \cdot u)\bigr)_i.
\]
This holds for every $1 \le i \le m$, so the two column vectors are equal.
\end{exercisesolution}

\begin{exercisesolution}[Proof of \cref{exc:continuous_subspace}]
We check the three conditions of \cref{def:linear_subspace} for $C(S) \subseteq \mathbb{R}^S$, with $S = [0,1]$. Each rests on a standard fact from analysis about limits, which is where this example steps outside algebra.

\textbf{(LSS1)} The constant function $0$ is continuous at every point of $[0,1]$, since $|0 - 0| = 0 < \varepsilon$ for every $\varepsilon > 0$ and every choice of $\delta$. So $0 \in C(S)$.

\textbf{(LSS2)} Let $f, g \in C(S)$ and let $x_0 \in [0,1]$. Given $\varepsilon > 0$, continuity of $f$ and of $g$ at $x_0$ supplies $\delta_1, \delta_2 > 0$ with $|f(x) - f(x_0)| < \varepsilon/2$ for $|x - x_0| < \delta_1$ and $|g(x) - g(x_0)| < \varepsilon/2$ for $|x - x_0| < \delta_2$. Put $\delta := \min(\delta_1, \delta_2)$. For $|x - x_0| < \delta$ the triangle inequality gives
\[
    \bigl|(f+g)(x) - (f+g)(x_0)\bigr| \le |f(x) - f(x_0)| + |g(x) - g(x_0)| < \tfrac{\varepsilon}{2} + \tfrac{\varepsilon}{2} = \varepsilon .
\]
So $f + g$ is continuous at $x_0$, and $x_0$ was arbitrary.

\textbf{(LSS3)} Let $f \in C(S)$ and $a \in \mathbb{R}$. If $a = 0$ then $a \cdot f$ is the zero function, covered by (LSS1). If $a \neq 0$, then given $\varepsilon > 0$ choose $\delta > 0$ with $|f(x) - f(x_0)| < \varepsilon / |a|$ for $|x - x_0| < \delta$; then $|a f(x) - a f(x_0)| = |a| \cdot |f(x) - f(x_0)| < \varepsilon$. So $a \cdot f$ is continuous.

Hence $C(S)$ is a linear subspace of $\mathbb{R}^S$. It is worth noticing that this is the first subspace in the book whose defining condition is not algebraic: nothing in the argument would work over an arbitrary field, because $|\cdot|$ and $\varepsilon$ have no meaning there.
\end{exercisesolution}

\begin{exercisesolution}[Proof of the Claim on Polynomials of Bounded Degree, on \cpageref{claim:bounded_degree_subspace}]
We use the criterion of \cref{lem:subspace_criterion}. The zero polynomial has degree $-\infty \le d$, so it lies in $K[x]_d$ and the set is non-empty. Now let $f, g \in K[x]_d$ and $a, b \in K$, and write
\[
    f = \sum_{k \ge 0} f_k x^k, \qquad g = \sum_{k \ge 0} g_k x^k,
\]
where $f_k = 0$ and $g_k = 0$ for every $k > d$. The coefficient of $x^k$ in $a f + b g$ is $a f_k + b g_k$, which vanishes for every $k > d$ as well. Hence $\deg(a f + b g) \le d$, so $a f + b g \in K[x]_d$.

Notice that the inequality $\deg(a f + b g) \le d$ can be strict even when $\deg f = \deg g = d$: over any field, $f = x^d$ and $g = -x^d$ give $f + g = 0$, of degree $-\infty$. That is exactly why the definition uses $\le d$ rather than $= d$, and it is the content of \cref{q:degree_exactly_d} \textbf{(a)}.
\end{exercisesolution}

\begin{exercisesolution}[Proof of the Claim on Matrix Spaces, on \cpageref{claim:matrix_space}]
This is the first solution above once more, with index set $I = \{1, \dots, m\} \times \{1, \dots, n\}$: a matrix is exactly a function assigning an element of $K$ to each pair $(i,j)$, and both operations are defined entry by entry. Every one of (V1)--(V8) is therefore an identity checked at each entry separately, where it becomes one of the field axioms. The zero vector is the matrix with $0$ in every entry, as claimed, and the additive inverse of $(a_{ij})$ is $(-a_{ij})$.

In the language of \cref{ex:function_space} this says that $\M_{m \times n}(K)$ is the function space $K^{I}$ for that index set. The identification is not merely formal: it is what makes $\dim \M_{m \times n}(K) = m \cdot n$ evident later, the standard basis being the matrix units, one for each index.
\end{exercisesolution}

\begin{answer}
Both parts of \cref{q:degree_exactly_d}.

\textbf{(a) No.} The set $\{f \in K[x] \mid \deg(f) = d\}$ is never a subspace, for any $d \in \mathbb{Z}_{\ge 0}$. The quickest reason is that it fails (LSS1): the zero polynomial has degree $-\infty \neq d$, so it is not in the set, and by \cref{lem:subspace_criterion} a subspace must contain $0$. Closure fails too, and more informatively: $x^d$ and $1 - x^d$ both have degree $d$, while their sum $1$ has degree $0$.

\textbf{(b) Yes.} Write $W$ for the set in question and let $\varphi \colon K[x]_5 \to K$ send $p = a_0 + a_1 x + \dots + a_5 x^5$ to $\sum_{i=0}^{5} a_i$. Then $W = \{p \in K[x]_5 \mid \varphi(p) = 0\}$, and we check the criterion of \cref{lem:subspace_criterion} directly. The zero polynomial has all coefficients $0$, whose sum is $0$, so $0 \in W$. And if $p, q \in W$ and $a, b \in K$, then the $i$-th coefficient of $a p + b q$ is $a p_i + b q_i$, so
\[
    \sum_{i=0}^{5} (a p_i + b q_i) = a \sum_{i=0}^{5} p_i + b \sum_{i=0}^{5} q_i = a \cdot 0 + b \cdot 0 = 0,
\]
and $a p + b q \in W$.

The contrast between the two parts is the point. A condition of the form \qt{some linear expression in the coefficients vanishes} always cuts out a subspace, because $0$ satisfies it and the expression respects linear combinations. A condition of the form \qt{some coefficient does \emph{not} vanish}, which is what $\deg f = d$ says about $f_d$, never does, for the same reason read backwards. Both halves of that observation are the substance of \cref{thm:solution_set_subspace} in a different costume, and $\varphi$ above is the first linear functional in the book, a notion ch. 18 takes up in earnest.
\end{answer}
